\section{A molecular inheritance realization}
\label{sec:source}

At $N$ sites a cell has $a$ activating, $r$ repressive and $s=N-a-r$ unmarked
sites. Following the reader--writer motif and the companion chronological
inheritance construction \citep{dodd2007,molecular}, the four chemical rates are
\begin{align}\label{eq:chem}
 (a,r)\to(a+1,r)&:\ s(.01+a/N),&
 (a,r)\to(a,r+1)&:\ s(.01+r/N),\\
 (a,r)\to(a-1,r)&:\ a\bigl(e+r/(2N)\bigr),&
 (a,r)\to(a,r-1)&:\ r\bigl(e+a/(2N)\bigr).\nonumber
\end{align}
Division removes the mother. The ordered daughters receive $(i,j)$ and
$(a-i,r-j)$ old marks with probability $2^{-(a+r)}\binom ai\binom rj$ and are
refilled to $N$ sites with unmarked material; there is no extra dilution term
and no artificial recovery interval. Division is $b=.1$, and death is $.01$ if
$a>r$ and $.3$ otherwise, so the states with $a>r$ are the \emph{protected}
ones. Intrinsic erasure is $e=.01+v$, so the eraser actuator has
$R=\partial_eQ_e$ and amplitude $v_{\max}=e_{\max}-.01$. The time unit is
synthetic and these rates are modelling assumptions, not measured
phenotype-switching rates.

For $N=2$ order the states as $(UU,UR,RR,AU,AR,AA)$. Their mark counts are,
respectively, $(0,0)$, $(0,1)$, $(0,2)$, $(1,0)$, $(1,1)$ and $(2,0)$.
A complete source specification is
\begin{align*}
 (Q_ef)_0&=(f_1-f_0)/50+(f_3-f_0)/50,\\
 (Q_ef)_1&=.51(f_2-f_1)+.01(f_4-f_1)+e(f_0-f_1),\\
 (Q_ef)_2&=2e(f_1-f_2),\\
 (Q_ef)_3&=.51(f_5-f_3)+.01(f_4-f_3)+e(f_0-f_3),\\
 (Q_ef)_4&=(e+.25)(f_1+f_3-2f_4),\\
 (Q_ef)_5&=2e(f_3-f_5),
\end{align*}
with the ordered daughter law
\[
 D(x,x)=\left(x_0^2,\ x_0x_1,\ \frac{x_0x_2+x_1^2}{2},\ x_0x_3,\
          \frac{x_0x_4+x_1x_3}{2},\ \frac{x_0x_5+x_3^2}{2}\right)^{\!T}.
\]
The one-daughter marginal follows from this law; the mean operator is its
linearization at $\one$, not an independent-daughter replacement. All exact
computations below use the pair law.

\subsection{Exact exposure certificates at two sites}
\label{sec:twosite}

Take
\begin{equation}\label{eq:q}
 q=(.94,.982,.987,.27,.735,.235)^T,\qquad c=67/73 .
\end{equation}
Exact rational calculation gives $\Phi_{.01}(q)\le0$ with minimum slack
$1/5000$. The derivative of $Q_e$ satisfies
\[
 Rq=\bigl(0,\ q_0-q_1,\ 2(q_1-q_2),\ q_0-q_3,\ q_1+q_3-2q_4,\
      2(q_3-q_5)\bigr)^T,
\]
and the slacks in $c(1-q)-Rq$ are exactly
\begin{equation}\label{eq:slacks}
 \left(\frac{201}{3650},\ \frac{534}{9125},\ \frac{1601}{73000},\ 0,\
       \frac{33669}{73000},\ \frac{9229}{14600}\right)^{\!T}.
\end{equation}
The certificate therefore holds for every nonnegative control increment for
which the bounded-action construction is made. The baseline continuation
satisfies $q^{\mathrm{off}}\le q$ by the supermartingale argument following
Theorem~\ref{thm:main}, which also covers any off-source with
$\Phi_{\mathrm{off}}(q)\le0$. For $n$ founders of type $AA$,
Theorem~\ref{thm:main} gives
\begin{equation}\label{eq:AA}
 r_\pi(T)\ \ge\ 1-\left[1-\frac{153}{200}e^{-67B/73}\right]^n .
\end{equation}
The zero slack at $AU$ makes $67/73$ the smallest admissible coefficient for
this particular $q$ and these six inequalities. It identifies a limiting
certificate constraint, not a dominant biological trajectory and not an optimal
intervention target.

\begin{remark}[How much is lost by the displayed certificate]
\label{rem:optcert}
The bound \eqref{eq:AA} improves on the horizon-independent form
$[1-(47/200)^n]e^{-67B/73}$ used in earlier drafts: writing $\alpha=e^{-cB}$,
the function $1-(1-\alpha\eta)^n$ is concave in $\alpha$ and lies above its
chord $\alpha[1-(1-\eta)^n]$, with equality at $n=1$, so the improvement is in
the founder dependence and not a new one-founder constant. One may also ask how
much is lost by fixing \eqref{eq:q}. Optimizing the certificate numerically and
then verifying the result in exact rational arithmetic gives
$\eta=.77849912\ldots$ and $c=.916533$ in place of $\eta=.765$ and
$c=.917808\ldots$, raising $B_{\mathrm{nec}}(1,.01)$ from $4.725705$ to
$4.751365$, an improvement of $0.54\%$. The vector is recorded in
the data file \texttt{site\_certificates.json}. The gap to the certified sufficient
exposure of Section~\ref{sec:pulse} is therefore not an artefact of the choice
\eqref{eq:q}.
\end{remark}

For the constant erasure $e=.3$, set
$w_*=(1419,1105,1000,8423,4305,10765)^T$. The exact residuals of
$-(21/200)w_*-A_{.3}w_*$ are
\[
 (201/200,\ 37/40,\ 11/20,\ 47/8,\ 99/40,\ 301/40)^T>0 ,
\]
so Proposition~\ref{prop:weight} applies with $\gamma=21/200$ and
$C=10.765$; equation \eqref{eq:upper} gives $10.765\,ne^{-.105t}$. With
$v_{\max}\ge.29$, an early $.29$ increment for
\begin{equation}\label{eq:sufficient2}
 \tau=\frac{200}{21}\log\frac{10.765n}{\delta},\qquad
 B_{\mathrm{suff}}=\frac{58}{21}\log\frac{10.765n}{\delta}
\end{equation}
achieves risk at most $\delta$ at every $T\ge\tau$ and after any subsequent
continuation without immigration. This certificate is for $e=.3$ exactly, not
for an unstated interval of uncertain erasure rates.

\begin{table}[t]
\centering
\caption{Two-site source, $AA$ founders, target risk $.01$, amplitude at least
$.29$. The sufficient horizon is part of the guarantee. Rounded values display
the exact logarithmic formulas \eqref{eq:necessary} and
\eqref{eq:sufficient2}; they are not rounded admission tests.}
\label{tab:budget}
\begin{tabular}{rrrr}
\toprule
$n$ & Necessary $B$ & Sufficient $B$ & Sufficient $T$\\
\midrule
1 & 4.725705 & 19.282156 & 66.490193 \\
10 & 7.229568 & 25.641677 & 88.419575 \\
100 & 9.737862 & 32.001198 & 110.348957 \\
1,000,000 & 19.772955 & 57.439281 & 198.066485 \\
\bottomrule
\end{tabular}

\end{table}

\begin{figure}[t]
\centering\includegraphics[width=.94\linewidth]{figures/budget.pdf}
\caption{Proved exposure floor and achievable weighted bound for the two-site
source, with free finite horizon. The gap is unclassified, not a numerically
estimated optimal frontier. Exposure is reaction rate integrated over synthetic
time.}
\label{fig:budget}
\end{figure}

\subsection{A certified finite-time pulse}
\label{sec:pulse}

For one $AA$ founder use $e=.3$ on $[0,40]$, then $e=.01$ on $(40,100]$, with
the same background killing throughout. The exposure is $B=.29\times40=11.6$.
The backward calculation is
$\mathcal F_{.3}^{40}\bigl(\mathcal F_{.01}^{60}(0)\bigr)$, where
$\mathcal F_e^t$ is the six-dimensional flow of $\dot u=\Phi_e(u)$. An
outward-rounded dyadic Taylor enclosure proves
\begin{equation}\label{eq:validated}
 \Prob_{AA}(|Z_{100}|>0)
 \ \le\ \frac{24892257049354995536083116261}
             {2535301200456458802993406410752}
 \ <\ 0.009818264214475\ <\ .01 .
\end{equation}
Appendix~\ref{app:validated} gives the error argument and arithmetic
specification. The sufficient exposure at this particular horizon is therefore
$11.6$, against $19.282156\ldots$ from \eqref{eq:sufficient2} and a necessary
$4.725704\ldots$ from \eqref{eq:necessary}; the certified bracket at $n=1$,
$T=100$, $v_{\max}=.29$, $\delta=.01$ is
\[
 4.725704\ldots\ \le\ B^*(100,.01)\ \le\ 11.6 ,
\]
a ratio of about $2.45$ rather than the $4.08$ obtained from the mean
certificate. Neither endpoint is asserted to be the exact optimum. For
independent founders under a \emph{deterministic} schedule, a single-founder
upper risk $\varepsilon$ gives population risk at most
$1-(1-\varepsilon)^n\le n\varepsilon$; this use of independence is legitimate
precisely because the schedule is predetermined. Equation
\eqref{eq:validated} alone therefore does not achieve $.01$ for ten or a
hundred founders.

\subsection{Preparation and withdrawal are part of the theorem}
\label{sec:withdrawal}

Three experiments must be distinguished: eraser withdrawal under continuing
background killing; withdrawal of both interventions into a specified
off-source; and a later rechallenge or regrowth-threshold assay. The last
requires its own stopped event and is not analyzed here. After eraser-only
withdrawal our default off-source is $e=.01$ with the original death vector. If
the off-source preserves $Q,b,D$ and reduces death coordinatewise then
$\Phi_{\mathrm{off}}(q)\le\PhiB(q)\le0$ and the exposure floor persists. For the
explicit full-withdrawal scenario with $b=.1$, $d=.01$ in every state, the total
count is an ordinary birth--death process with extinction vector
$q^{\mathrm{off}}=.1\,\one$, the smaller root of $.01(1-x)+.1(x^2-x)=0$. This is
a stipulated full-withdrawal model, not a fitted biological law.

\begin{figure}[t]
\centering\includegraphics[width=\linewidth]{figures/policy.pdf}
\caption{Left: $AA$ survival by time $100$ for equal-exposure early, late and
spread eraser schedules (lines: numerical diagnostics). The marked upper bound
at $B=11.6$ is rigorous. Right: numerical eventual-survival comparisons after
eraser withdrawal at $40$, keeping background killing until $100$ (then
baseline or full withdrawal), or removing all killing immediately at $40$.
Every plotted continuation is stated; no curve is a clinical relapse estimate.}
\label{fig:policy}
\end{figure}

Erasure is not universally beneficial at a fixed short horizon. From an $RR$
founder the analytic finite-time extinction solution satisfies
\begin{equation}\label{eq:taylor}
 u_{RR}^{.3}(t;0)-u_{RR}^{.01}(t;0)
 =-\frac{49619}{600000000}\,t^4+O(t^5),
\end{equation}
so stronger erasure \emph{lowers} extinction probability for all sufficiently
small $t>0$. This is a source-level failure of monotonicity: erasure also
removes repressive marks and permits access to the protected active states, and
the source assigns lower death to $a>r$. It is a rigorous counterexample to
universal monotonicity, not evidence that an epigenetic drug rescues tumours,
and it does not imply a large experimental effect: at $t=1$ the numerical
difference is about $4.1\times10^{-5}$. The exact calculation is in
Appendix~\ref{app:taylor}; finite-time numerical differences are labelled
separately. Preparation, budget, duration and endpoint all matter, and an
early/late ordering seen for $AA$ need not persist for $RR$.

\subsection{Slower protected division}
\label{sec:variation}

To challenge the demographic assumptions without rescaling the clock, replace
the division rate in states $a>r$ by a common unknown fixed $b_P\in[.08,.1]$,
keeping all death and chemical rates fixed. An independently checked rational
supersolution and positive weight give, uniformly on that interval,
\[
 q_{AA}=.34796677,\qquad c=0.9163811468\ldots,\qquad
 \gamma=.1004,\qquad C=10.764704 .
\]
The data file \texttt{source\_certificates.json} records the exact value of
$c$, all coordinates and the endpoint residuals. Each inequality is affine in
$b_P$, so
exact endpoint checks prove the whole interval. This is a scenario family, not a
confidence interval inferred from experiments, and a $20\%$ reduction in
protected division is a modest perturbation; it does not demonstrate robustness
to deep dormancy or to orders-of-magnitude slower persister division.
